\specialchapter{INTRODUCTION}

In Chapters III to V of this Book, we consider in detail the algebra of endomorphisms of a real or complex topological vector space. In the case of finite-dimensional spaces, the principal properties of the elements of this algebra, taken individually, follow from the Jordan decomposition (A, VII, p. 42, § 5, no. 9). Here we study their generalization to infinite-dimensional spaces, which is however more delicate.

In Chapter III, we consider, in a very general framework, compact linear maps between topological vector spaces. These are the ones whose spectral properties most closely resemble the finite-dimensional case. For example, in the particular case of Banach spaces, the nonzero elements of the spectrum of a compact endomorphism are generalized eigenvalues of finite spectral multiplicity.

In the context of Hilbert spaces, a manageable classification and structure emerge. This is the object of Chapter IV. It begins with the study of compact operators in Hilbert spaces: their description by singular values, as well as their diagonalization in the case of a normal endomorphism, are entirely analogous to the properties known for a finite-dimensional space. More generally, the spectral theorem provides simple models of normal endomorphisms of Hilbert spaces.

Certain important applications, concerning for example differential operators, but also quantum mechanics, require a more subtle theory: that of linear maps --- not necessarily continuous --- defined on a dense subspace of a Hilbert space, which we call partial operators. In this framework, the very definition of a self-adjoint or normal operator demands great precision if one wishes to avoid gross blunders. Once these definitions are acquired, a classification similar to that of continuous endomorphisms emerges; the functional calculus for universally measurable maps, bounded or not, provides close links between continuous normal endomorphisms and normal partial operators.

In passing, we shall have mentioned normal endomorphisms with trace, and in particular integral endomorphisms with kernel in a Hilbert space \(L^{2}(X,\mu)\) . One knows the crucial importance, in the theory of automorphic forms and representations, of the formula establishing that the trace of such an endomorphism, when it exists, is the integral of the kernel over the diagonal. We also devote a paragraph to the fundamental definitions and properties of distributions and tempered distributions in \(R^{n}\) . These notions provide not only the most convenient presentation of the fundamental principles of Euclidean Fourier theory (cf. II, § 1), but they also allow one to identify the domains of many important differential operators.

Finally, in Chapter V, we relate the study of endomorphisms of Hilbert spaces to that of unitary representations of groups. After some elementary general considerations, among which Schur's lemma is the most important statement, we study the case of locally compact groups. We then consider functions of positive type, and present their first applications. One of these is the Gelfand-Naimark-Segal theorem, which asserts that every star algebra is isomorphic to a closed subalgebra of the algebra of endomorphisms of a Hilbert space. Others concern fundamental properties of the space of irreducible unitary representations of a locally compact group, in particular the Gelfand-Raikov theorem, which establishes that these representations separate points. In the last paragraph of this chapter, we prove the fundamental results concerning representations of compact groups, which served us in Chapter IX of Groups and Lie algebras.

A. BOREL, Representations of locally compact groups, Lecture Notes in Mathematics 276, Springer, 1972.

M. REED and B. SIMON, Methods of Mathematical Physics, Academic Press, Vol. 1, 1980 and Vol. 2, 1975.

K. SCHMÜDGEN, Unbounded Self-adjoint Operators on Hilbert Space, Graduate Texts in Mathematics 265, Springer, 2012.

In the preceding chapters and books, the following references to “TS, to appear”, must be modified as follows:

EVT, V, p. 35, line 1. Instead of “TS, IV, § 6”, read “(TS, IV, p. 160, remark)”.

INT, IX, p. 94, § 6, no. 12. Instead of “One can prove a converse \ldots{} bounded”, read “The converse of prop. 11 is a particular case of TS, V, p. 455, th. 5.”

LIE, IX, p. 61, § 6, no. 4, line 6. Instead of “TS, to appear”, read “TS, III, p. 36, prop. 9”.

LIE, IX, p. 63, § 6, no. 4, line 6. Instead of “TS, to appear”, read “TS, III, p. 36, rem.”.

LIE, IX, p. 66, § 7, note 1. Instead of “to a chapter to appear”, read “to paragraph 4 of Chapter V”.

LIE, IX, p. 69, § 7, no. 2. Instead of “then according to TS”, read “then (TS, V, p. 464, prop. 6)”.

LIE, IX, p. 71, § 7, no. 3. Instead of “which, according to TS, is an isomorphism”, read “which, according to TS, V, p. 469, cor. 2, is an isomorphism.”

LIE, IX, p. 72, § 7, no. 4. Instead of “(TS)”, read “(TS, V, p. 459, cor. 3)”.

LIE, IX, p. 74, § 7, no. 4. After “of the space of central representative functions on G”, add “(TS, V, p. 469, cor. 2, a)”; after “of the space ZL²(G) of central square-integrable functions on G”, add “(TS, V, p. 468, prop. 9).”

LIE, IX, p. 78, § 8, no. 1. Instead of “In this number, one recalls definitions and results from TS”, read “One recalls definitions and results from TS, V, p. 386, no. 8; p. 392, no. 10; p. 470, no. 7.”

LIE, IX, p. 84, § 8, no. 3. Instead of “Recall (TS) the formulas”, read “One has (TS, V, p. 467, formulas (1) and (2)) the formulas”.

LIE, IX, p. 89, § 9, no. 1. Instead of “TS, III, §2, no. 7, prop. 16”, read “TS, III, p. 67, prop. 10”.

LIE, IX, p. 90, § 9, no. 1. Instead of “TS, III, §2, no. 6”, read “TS, III, p. 61, no. 4”.

LIE, IX, p. 91, § 9, no. 2. Instead of “TS, to appear”, read “TS, III, p. 36, prop. 9 and TS, V, p. 464, prop. 7”.

LIE, IX, p. 92, note 1. Instead of “(TS, to appear)”, read “(TS, V, p. 379, def. 3)”.

LIE, IX, p. 99, appendix 1, no. 1. Instead of “(TS, to appear)”, read “(TS, V, p. 464, remark 4)”.

LIE, IX, p. 128, exercise 3 (b) of §9. Instead of “(TS, to appear)”, read “(Use the Peter--Weyl theorem, TS, V, p. 462, th. 1)”.

TS, I, p. 103, example 6. Instead of “V, to appear”, read “V, p. 442, th. 2.”

TS, I, p. 173, exerc. 5, d). Instead of “III, to appear”, read “III, p. 82, § 6.”

TS, II, p. 237, no. 9. Instead of “IV, to appear”, read “IV, p. 196, § 3”.

TS, II, p. 251, no. 1. Instead of “V, to appear”, read “V, p. 425, prop. 10”.

